% ============================================================
% 本文件由 papex-build.py 自动生成，请勿手工编辑。
% 修改请在 papex.json 中进行，随后重新构建。
% ============================================================

学术文献之间的引用关系构成了一张庞大的知识网络。准确预测一篇新论文可能引用的文献，对于文献推荐、科研趋势分析与学术检索具有重要意义。本文提出一种面向中文文献的图神经网络引用预测框架 Papex-GNN：首先利用预训练语言模型对中文标题与摘要进行语义编码，随后将其作为节点特征接入引文图，通过异构图注意力网络同时建模“主题相似”与“作者协作”两类边。在自建的中文论文数据集上，本文方法相较基于 TF-IDF 与协同过滤的基线，在 Recall@10 指标上提升了 12.4 个百分点。消融实验进一步表明，语义编码与结构建模两个模块均对最终效果有显著贡献。
